\section{Find the Nth Root of a Number using Binary Search}
\label{sec:bs-nth-root}

\problemstatement{Given two positive integers $n$ and $m$, find $x$ such
that $x^n = m$, where $x$ is a positive integer. If no such integer $x$
exists, report $-1$.}

\begin{patternbox}
This is the direct generalization of Section~\ref{sec:bs-sqrt} from
``square root'' to ``$n$-th root.'' The keyword pattern is identical:
\emph{``find the value $x$ such that $x$ raised to some power equals (or
is bounded by) a target.''} The only change is the predicate's exponent,
from $2$ to a general $n$ -- everything else about the binary-search-on-
the-answer template carries over unchanged.
\end{patternbox}

\subsection*{Understanding the problem carefully}

Define $P(\textit{candidate}) = (\textit{candidate}^n \le m)$ for
candidates in $[1, m]$. Just as squaring is monotonic for non-negative
integers, raising to any fixed positive power $n$ is monotonic too:
$x \le y \Rightarrow x^n \le y^n$. So $P$ is \texttt{true} for a prefix of
candidates and \texttt{false} afterward, with a single crossover -- the
largest candidate satisfying $P$ is a perfect $n$-th root only if
$\textit{candidate}^n$ exactly equals $m$ at that point; otherwise, no
exact integer $n$-th root exists.

\begin{ideabox}[Candidate Space and Predicate]
The candidate space is integers $[1, m]$. The predicate is
$P(\textit{candidate}) = (\textit{candidate}^n \le m)$. Track the largest
candidate for which $P$ holds; at the end, check whether that candidate,
raised to the $n$-th power, exactly equals $m$.
\end{ideabox}

\subsection*{Why we must verify exactness at the end}

Unlike the square-root problem, which explicitly asked for the
\emph{floor} of the root (so any final value of \textit{ans} is a valid
answer), this problem demands an \emph{exact} integer root, or $-1$ if
none exists. So the binary search itself is unchanged -- it always finds
the largest candidate whose $n$-th power does not exceed $m$ -- but we add
one final check: if that candidate's $n$-th power is not exactly $m$,
there is no integer solution, and we must report $-1$ rather than the
nearest approximation.

\subsection*{Algorithm}

\begin{enumerate}
  \item Initialize $\textit{lo} \gets 1$, $\textit{hi} \gets m$.
  \item While $\textit{lo} \le \textit{hi}$:
  \begin{enumerate}
    \item $\textit{mid} \gets \textit{lo}+(\textit{hi}-\textit{lo})/2$.
    \item Compute $p \gets \textit{mid}^n$ (stop early/clamp if it exceeds
          $m$, to avoid overflow for large $n$).
    \item If $p = m$: return \textit{mid}.
    \item If $p < m$: $\textit{lo} \gets \textit{mid}+1$.
    \item Else: $\textit{hi} \gets \textit{mid}-1$.
  \end{enumerate}
  \item If the loop ends without returning, return $-1$.
\end{enumerate}

\begin{algorithm}[H]
\caption{Nth Root}
\begin{algorithmic}[1]
\Procedure{NthRoot}{$n, m$}
  \State $\textit{lo} \gets 1$, $\textit{hi} \gets m$
  \While{$\textit{lo} \le \textit{hi}$}
    \State $\textit{mid} \gets \textit{lo}+(\textit{hi}-\textit{lo})/2$
    \State $p \gets \textit{mid}^n$ \Comment{computed with early-exit on overflow}
    \If{$p = m$} \Return \textit{mid} \EndIf
    \If{$p < m$}
      \State $\textit{lo} \gets \textit{mid} + 1$
    \Else
      \State $\textit{hi} \gets \textit{mid} - 1$
    \EndIf
  \EndWhile
  \State \Return $-1$
\EndProcedure
\end{algorithmic}
\end{algorithm}

\subsection*{Dry run}

\dryrun{Take $n = 3$, $m = 27$ (find the cube root).}

\begin{itemize}
  \item $\textit{lo}=1,\textit{hi}=27$: $\textit{mid}=14$, $14^3=2744 \le
        27$? No, $p > m$. $\textit{hi}=13$.
  \item $\textit{lo}=1,\textit{hi}=13$: $\textit{mid}=7$, $7^3=343 > 27$.
        $\textit{hi}=6$.
  \item $\textit{lo}=1,\textit{hi}=6$: $\textit{mid}=3$, $3^3=27=27$.
        Return $3$.
\end{itemize}

The answer is $3$, since $3^3 = 27$.

Now take $n=4$, $m=20$. No integer fourth root of $20$ exists ($2^4=16$,
$3^4=81$), so the binary search will narrow down without ever finding
$p=m$, and the final answer will correctly be $-1$.

\subsection*{C++ implementation}

\begin{lstlisting}
#include <bits/stdc++.h>
using namespace std;

// Computes base^exp, capped to indicate ">m" without overflow.
long long power(long long base, int exp, long long cap) {
    long long result = 1;
    for (int i = 0; i < exp; i++) {
        result *= base;
        if (result > cap) return cap + 1; // overflow guard
    }
    return result;
}

int nthRoot(int n, int m) {
    int lo = 1, hi = m;

    while (lo <= hi) {
        int mid = lo + (hi - lo) / 2;
        long long p = power(mid, n, m);

        if (p == m) return mid;
        else if (p < m) lo = mid + 1;
        else hi = mid - 1;
    }
    return -1;
}

int main() {
    cout << "Cube root of 27: " << nthRoot(3, 27) << endl;
    cout << "4th root of 20: " << nthRoot(4, 20) << endl;
    return 0;
}
\end{lstlisting}

\begin{complexitybox}
\textbf{Time Complexity.} The candidate range $[1, m]$ halves each
iteration, and each predicate evaluation costs $O(n)$ to compute the
power (or $O(\log n)$ with fast exponentiation), giving an overall time
complexity of
\[
  O(n \log m)
\]
(or $O(\log n \log m)$ with fast exponentiation).
\textbf{Space Complexity.} $O(1)$ auxiliary space.
\end{complexitybox}

\begin{takeawaybox}
Generalizing a binary-search-on-the-answer problem from a fixed exponent
(square root, exponent $2$) to a variable exponent ($n$-th root) changes
nothing about the search itself -- only the cost of evaluating the
predicate changes, since computing $\textit{candidate}^n$ is more expensive
than $\textit{candidate}^2$. Always separate ``how do I search the
candidate space'' from ``how expensive is one feasibility check,'' since
they contribute independently to the total complexity.
\end{takeawaybox}
